%% Khối tiêu đề, tóm tắt và từ khoá (bản tiếng Việt). Được access_vi.tex \input ngay sau \begin{document}.
%% TODO: danh sách tác giả, e-mail và DOI trước khi nộp.
\history{Ngày công bố xxxx 00, 0000, ngày của phiên bản hiện tại xxxx 00, 0000.}
\doi{10.1109/ACCESS.2026.0000000}

\title{Từ phương trình đến điện cực: kiểm định có đăng ký trước mô hình Bộ não như một giao diện truy vấn của thiền trên hai bộ dữ liệu EEG mở}
% TODO: xác nhận danh sách tác giả và đơn vị trước khi nộp (cùng khối như bài đồng hành).
\author{\uppercase{Dung Vu}\authorrefmark{1}}
\address[1]{Trường Đại học Phenikaa, Hà Nội 12116, Việt Nam (e-mail: [e-mail cơ quan])}
\tfootnote{Nghiên cứu được thực hiện tại Trường Đại học Phenikaa, Hà Nội. Mọi tính toán chạy trên máy tính cá nhân; không nhận tài trợ bên ngoài.}

\markboth
{Vu: Từ phương trình đến điện cực: kiểm định mô hình BQI của thiền trên EEG}
{Vu: Từ phương trình đến điện cực: kiểm định mô hình BQI của thiền trên EEG}

\corresp{Tác giả liên hệ: Dung Vu (e-mail: [e-mail cơ quan]).}

\begin{abstract}
Các mô hình động lực nối hoạt động của mạng chế độ mặc định, mức tích hợp thông tin và chất lượng thiền do người tập tự báo cáo ngày càng được đề xuất nhiều, và một phân tích thiết kế dựa trên thông tin Fisher chính xác của một mô hình không gian trạng thái tuyến tính--Gauss gần đây đã dự đoán mô hình như vậy có thể được ước lượng chính xác đến đâu từ điện não đồ (EEG) lấy mẫu trải nghiệm. Các dự đoán đó có đúng trên bản ghi thật hay không thì chưa được kiểm tra. Chúng tôi fit mô hình lên hai bộ dữ liệu EEG lấy mẫu trải nghiệm mở của hai phòng thí nghiệm khác nhau, ds001787 (24 người tập, 40 phiên, 983 lần hỏi tự chấm) và OSF 3uszv (58 người, 58 phiên, 2\,258 lần hỏi), theo hai bản đăng ký trước được khoá bằng mã băm và bốn cổng kiểm tra tự động cho mỗi bộ. Trên cả hai bộ, sai số tương đối của hằng số tốc độ của mạng chế độ mặc định được dự đoán trước khi khoá gần với sai số chuẩn Fisher tại ước lượng và với sai số bootstrap tham số (9,1; 9,9 và 9,7 phần trăm trên ds001787; 6,7; 8,5 và 8,8 phần trăm trên OSF 3uszv), trong khi sai số theo thông tin quan sát của kênh công suất alpha vượt giá trị Fisher, một dấu hiệu sai đặc tả không xuất hiện trên dữ liệu mô phỏng từ fit. Không đại diện EEG nào cải thiện dự đoán ngoài mẫu của điểm tự chấm so với các mô hình rỗng trên bộ nào. Bản thân các hằng số tốc độ không lặp lại: hằng số thời gian của mạng chế độ mặc định là 0,76 phút trên ds001787 và 7,1 phút trên OSF 3uszv, và một mô phỏng cho thấy fit Ornstein--Uhlenbeck trả về thang thời gian trội của bất kỳ đặc trưng nào được đưa vào. Các phân tích thiết kế kiểu này đáng tin về độ chính xác mà chúng hứa hẹn, nhưng không đáng tin về ý nghĩa của các tham số ước lượng từ đại diện chưa hiệu chuẩn.
\end{abstract}

\begin{keywords}
Bộ lọc Kalman, cận Cramér--Rao, đăng ký trước, điện não đồ, lấy mẫu trải nghiệm, mạng chế độ mặc định, mô hình không gian trạng thái, sai đặc tả mô hình, thiền, thông tin Fisher.
\end{keywords}

\titlepgskip=-21pt

\maketitle
